\section{Discussion}\label{sec:discussion}
Internal accuracy of \RTestAcc\% did not predict external behaviour: the same model labels only
\RExtSpecK/\RExtSpecN{} normal Montgomery scans Normal and names none of the \RExtTBN{} TB scans
Tuberculosis. The internal split shares sources, acquisition styles and labelling practices with training,
even when patients are separated; an external cohort breaks that link
\cite{zech2018variable,degrave2021ai,geirhos2020shortcut}.

\emph{Source confounding.} Class and source overlap in the merged data. In run 1 the external scans were
frequently called COVID-19, a class supplied only by web-collected sources. COVID-19 calls fell from \RRunOneExtCovidCalls{} to
\RExtCovidCalls{} of \RNExternal{} with build 2, yet specificity stayed low and the errors moved to Pneumonia and Nodule/Mass.
Build 2 changed several sources at once (TBX11K, the COVID-19 Radiography Database and further NIH archives), so
this is consistent with, but not proof of, the added COVID-19 source removing a source cue. Our evidence is
indirect, and we do not claim that the model uses source cues.

\emph{Limitations.}
(i) One training run per model (seed \RSeed), CPU-trained and time-boxed; run 2's selection advantage is
largest on sources that only it saw. (ii) NIH labels are mined from reports, not read from images.
(iii) Patient identity is unknown for TBX11K and COVID-19 Radiography images, so residual patient overlap
across splits is possible. (iv) Kermany images are paediatric, the other sources adult.
(v) Pleural Effusion and Nodule/Mass come from NIH only. (vi) The single-label design excludes
multi-finding images. (vii) The external cohort is small and binary, so intervals are wide.
(viii) Data are retrospective and public; no clinical validity is claimed.
Further analyses, including seed variability, are ongoing and are not reported here.
